\begingroup%
\definecolor{ksTeal}{HTML}{27877F}%
\definecolor{ksAmber}{HTML}{C58B28}%
\definecolor{ksInk}{HTML}{243447}%
\definecolor{ksPaleTeal}{HTML}{EFF8F5}%
\definecolor{ksPaleAmber}{HTML}{FFF8E8}%
\begin{tikzpicture}[
  x=1mm, y=1mm,
  card/.style={draw, fill=white, rounded corners=2mm, line width=.7pt,
    text=ksInk, align=flush center, font=\fontsize{9.5}{12.5}\selectfont,
    minimum width=66mm, minimum height=15mm, inner xsep=2mm, inner ysep=1.7mm,
    outer sep=.35pt},
  frame/.style={draw, rounded corners=3mm, line width=.7pt, inner sep=3mm,
    outer sep=.35pt},
  heading/.style={text=ksInk, font=\bfseries\fontsize{10}{12.5}\selectfont,
    inner sep=0pt, outer sep=0pt},
  flow/.style={-{Stealth[length=2.25mm,width=1.65mm]}, line width=.7pt},
  edge text/.style={text=ksInk, font=\fontsize{9}{12.5}\selectfont, inner sep=1mm}
]
  \node[heading] (human-title) at (37,0)
    {人类 HPC 算子专家：将算子视为 Design 搜索问题};
  \node[card, draw=ksTeal] (h1) at (0,-14)
    {硬件规格与 Microbench\\
    测算各单元 Cycle 延迟与吞吐};
  \node[card, draw=ksTeal] (h2) at (74,-14)
    {锁定核心物理瓶颈\\
    如 B200 FA4 的 SFU 节拍瓶颈};
  \node[card, draw=ksTeal] (h3) at (74,-33)
    {纯粹的高层架构设计\\
    数学改写 · Warp 专职分工\\
    SMEM/TMEM 配平 · 异步掩盖};
  \node[card, draw=ksTeal] (h4) at (0,-33)
    {工程编码落地\\
    按既定蓝图翻译为底层代码};

  \node[heading] (agent-title) at (37,-60)
    {当前编码 Agent：将算子降格为短视的 Code 工程任务};
  \node[card, draw=ksAmber] (a1) at (0,-74)
    {长程奖励极度稀疏\\
    唯一目标是尽快跑通交差};
  \node[card, draw=ksAmber] (a2) at (74,-74)
    {跳过瓶颈与架构推演\\
    上来立刻埋头写底层代码};
  \node[card, draw=ksAmber] (a3) at (74,-93)
    {缺乏承载设计的结构化语言\\
    仅靠散漫 Markdown 泛泛记述};
  \node[card, draw=ksAmber] (a4) at (0,-93)
    {陷入恶性循环\\
    既无设计 Taste，又无演化载体};

  \begin{scope}[on background layer]
    \node[frame, draw=ksTeal, fill=ksPaleTeal,
      fit=(human-title)(h1)(h2)(h3)(h4)] (human) {};
    \node[frame, draw=ksAmber, fill=ksPaleAmber,
      fit=(agent-title)(a1)(a2)(a3)(a4)] (agent) {};
  \end{scope}

  \draw[flow, ksTeal] (h1) -- (h2);
  \draw[flow, ksTeal] (h2) -- (h3);
  \draw[flow, ksTeal] (h3) -- (h4);
  \draw[flow, ksAmber] (a1) -- (a2);
  \draw[flow, ksAmber] (a2) -- (a3);
  \draw[flow, ksAmber] (a3) -- (a4);
  \draw[flow, ksInk, dashed] (human.south) --
    node[edge text, right=1mm] {核心差异：Design 先行 vs. 盲目 Coding} (agent.north);
\end{tikzpicture}%
\endgroup%
